% TOP_PROBLEM: {"id": 30006749, "title": "Indistinguishability Obfuscation from Constant-Degree Homomorphic Assumptions", "category_id": 15, "category": {"id": 15, "name": "computer_science", "display_name": "Computer Science", "description": "Computational complexity, algorithms, and theoretical CS.", "slug": "computer-science", "order_index": 15, "created_at": "2026-07-31T15:26:25.671Z"}, "status": "open"}
% FIRST_POSTED: null
% !TeX program = lualatex
% Compile twice: lualatex 30006749.tex
% All references and prose are embedded; no BibTeX database or external inputs.
\documentclass[11pt,a4paper]{article}
\usepackage[margin=24mm,headheight=14pt]{geometry}
\usepackage{fontspec}
\setmainfont{Latin Modern Roman}
\setsansfont{Latin Modern Sans}
\setmonofont{Latin Modern Mono}
\usepackage{amsmath,amssymb,mathtools}
\usepackage{unicode-math}
\setmathfont{Latin Modern Math}
\usepackage{microtype}
\usepackage{enumitem,xurl,needspace}
\usepackage[dvipsnames]{xcolor}
\usepackage{hyperref}
\hypersetup{unicode=true,colorlinks=true,linkcolor=MidnightBlue,urlcolor=MidnightBlue,
 pdftitle={Research attempt: Problem 30006749},pdfauthor={Research notebook},bookmarksnumbered=true}
\usepackage{bookmark}
\usepackage{tocloft}
\setlength{\cftsecnumwidth}{3em}
\cftsetpnumwidth{2.5em}
\cftsetrmarg{4em}
\usepackage{titlesec}
\titleformat{\section}{\Large\bfseries\raggedright}{\thesection}{0.7em}{}
\usepackage{fancyhdr}
\pagestyle{fancy}\fancyhf{}
\fancyhead[L]{\small\sffamily Open Problems Math | Research notebook}
\fancyhead[R]{\small Problem 30006749 | 27 September 2026}
\fancyfoot[C]{\thepage}
\setlength{\parindent}{0pt}
\setlength{\parskip}{3pt plus 1pt}
\setlength{\emergencystretch}{3em}
\setcounter{tocdepth}{1}
\setcounter{secnumdepth}{1}
\setlist[itemize]{leftmargin=3em,itemsep=5pt,topsep=4pt,parsep=0pt}
\allowdisplaybreaks
\newcommand{\N}{\mathbb{N}}
\newcommand{\Z}{\mathbb{Z}}
\newcommand{\Q}{\mathbb{Q}}
\newcommand{\R}{\mathbb{R}}
\newcommand{\C}{\mathbb{C}}
\newcommand{\F}{\mathbb{F}}
\newcommand{\A}{\mathbb{A}}
\newcommand{\PP}{\mathbb P}
\newcommand{\E}{\mathbb{E}}
\newcommand{\eps}{\varepsilon}
\newcommand{\dd}{\,\mathrm d}
\newcommand{\sref}[2]{\hyperlink{source:#1:#2}{[S#2]}}
\newcommand{\eref}[2]{\hyperlink{extra:#1:#2}{[E#2]}}
% Turn the next switch off for a shorter reading copy. The quotations preserve
% every exactTarget field, including qualifications omitted from a short summary.
\newif\ifshowcatalogue
\showcataloguetrue
\newcommand{\cataloguescope}[1]{\ifshowcatalogue
 \paragraph{Exact scope in the supplied catalogue.}
 \begin{quote}\small #1\end{quote}\fi}
\usepackage{etoolbox}
\AtBeginEnvironment{itemize}{\small}
\numberwithin{equation}{section}
% Standalone export. Original research content preserved.
\begin{document}
\setcounter{section}{0}
% ================================================================
% problemId: 30006749
% Canonical problem ID: 30006749
% exactTarget SHA-256: ad29da5c92ca52381164224b9d1cd216c9a57201cef570bb86a55d38da86d33e
\clearpage
\section*{\texorpdfstring{Indistinguishability Obfuscation from Constant-Degree Homomorphic Assumptions}{Indistinguishability Obfuscation from Constant-Degree Homomorphic Assumptions}}\label{30006749}
\subsection{Definitions and mathematical statement}
An indistinguishability obfuscator is a randomized efficient transformation $C\mapsto\mathcal O(C)$ preserving the circuit's function and producing only polynomial overhead. For same-size equivalent circuits $C_0,C_1$, it requires
\[
 \mathcal O(C_0)\approx_c\mathcal O(C_1),
\]
meaning every probabilistic polynomial-time distinguisher has negligible advantage, in the usual security-parameter experiment. The goal is such a construction from an explicitly stated combination of standard polynomial-security assumptions of the indicated type, without subexponential security losses. ``Standard assumptions'' is not itself a formal predicate: a complete theorem must name the exact LWE/LPN/pairing variants, parameters, and reduction loss. No unrestricted high-degree multilinear-map assumption is silently added.
\par\smallskip\textit{Formulation and concept references:} \sref{30006749}{1}.
\cataloguescope{Construct indistinguishability obfuscation based purely on standard polynomial-time assumptions (such as standard LWE, LPN, and degree-2 multilinear pairings) without sub-exponential loss.}
\subsection{Short English statement}
Can general indistinguishability obfuscation be obtained from the stated kinds of conventional assumptions without requiring subexponential security?
% BEGIN RESEARCH ADDITION 142
\subsection{Research attempt}

\paragraph{Current frontier and assumption audit.}
The full version of the supplied STOC construction, Theorem 5.5, assumes
subexponential hardness of SXDH, a specified large-modulus LWE problem,
sparse-error LPN over $\Z_p$, and a polynomial-stretch Boolean
$\mathsf{NC}^0$ PRG. Its polynomial-security conclusion for collusion-resistant
functional encryption is not general iO. \sref{30006749}{1}\eref{30006749}{1}
Proof-sensitive obfuscation avoids input-by-input security loss for circuits
with short proofs of equivalence; the 2025 quasi-linear construction still
retains that hypothesis. \eref{30006749}{2}\eref{30006749}{3}
An ideal pseudorandom-oracle construction is in a different model, not an
instantiation from the requested standard assumptions. \eref{30006749}{4}

\paragraph{Attack route.}
Replace exponentially many inputwise hybrids by a short sequence of
semantics-preserving circuit rewrites, perhaps through proof compression.
We investigate whether short paths suffice for security, whether every
equivalent pair has one, and whether a reduction can actually access it.
None of the latter two assertions follows from semantic equivalence alone.

\paragraph{Attempt: exact accounting of hybrid loss.}
For distributions $H_0,\ldots,H_h$ and a test $D$, let
$d_i=\Pr[D(H_i)=1]$. Telescoping gives
\begin{equation}\label{30006749:hybrid}
 |d_h-d_0|\le\sum_{i=1}^h|d_i-d_{i-1}|.
\end{equation}
If an edge experiment samples $I$ uniformly and supplies $H_{I-1}$ or $H_I$,
using the same $D$ has signed gap $(d_h-d_0)/h$. An exponential range of
indices needs only $\log_2h$ random bits, so an efficiently samplable indexed
hybrid can save time; it does not remove the advantage loss.

For instance $2^{-\sqrt\lambda}$ is negligible but multiplying it by
$2^\lambda$ does not give a vanishing bound. More generally, for any fixed
polynomial rescaling $\kappa=\lambda^b$ and any $a>0$, the negligible function
\[
 \epsilon(\kappa)=2^{-(\log_2\kappa)^2}
\]
satisfies $2^{\lambda^a}\epsilon(\lambda^b)\to\infty$. Therefore no fixed
polynomial parameter inflation makes an exponential hybrid loss a logical
consequence of arbitrary polynomial security. This is an obstruction to
that inference, not to every possible cryptographic reduction.

Suppose instead a uniform procedure supplies a path of at most polynomially
many rewrites between a given equivalent padded circuit pair of size $s$.
Suppose each indexed rewrite has a polynomial-time simulation and a
polynomial-loss reduction to a fixed finite collection of named assumptions,
uniformly in the path. For $s\le\operatorname{poly}(\lambda)$, polynomial
security and \eqref{30006749:hybrid} give negligible endpoint advantage.
This is a valid composition statement; neither an edge-secure encoding nor
a universal path procedure is constructed here. Merely verifying a supplied
path is not the same as producing it.

\paragraph{Why universal proof compression is a deep extra demand.}
Let $\mathcal R$ be any sound, polynomial-time decidable circuit rewrite
relation. Auxiliary wires and expansions are allowed but counted in the
full certificate encoding. If one polynomial bounds the total encoding
length of a path between \emph{every} equivalent circuit pair, then
\begin{equation}\label{30006749:collapse}
 \mathsf{NP}=\mathsf{coNP}.
\end{equation}
Given a formula $\varphi$, compare its circuit with constant one, padding
as needed. Guess a path of the promised polynomial size and check it.
Soundness excludes nontautologies; the universal bound certifies every
tautology. Thus $\mathsf{TAUT}\in\mathsf{NP}$, proving
\eqref{30006749:collapse}. This implication needs only existence of short paths,
not a path-finding algorithm. A total polynomial-time path finder with this
sound verifier would even decide equivalence and imply
$\mathsf P=\mathsf{NP}$. Neither implication is an implication of iO itself;
it is attached to this particular proposed method.

Similarly, a polynomial-time canonical representation satisfying
\[
 C_0\equiv C_1\quad\Longleftrightarrow\quad
 \operatorname{can}(C_0)=\operatorname{can}(C_1)
\]
would decide tautology by comparison with $\operatorname{can}(1)$ and imply
$\mathsf P=\mathsf{NP}$. A truth table supplies perfect semantic
canonicalization but uses $2^n$ output bits. It does not meet polynomial
overhead. Randomness and computational indistinguishability cannot simply
be replaced by canonicalization.

\paragraph{A restricted bilinear-encoding test.}
Consider a generic asymmetric bilinear model whose source exponents are
affine linear in hidden variables $x_1,\ldots,x_m$. Permit exponent addition,
multiplication by known scalars independent of those hidden variables, and
pairing two source elements; prohibit a target-to-source map. Induction
shows that source exponents remain affine and target exponents have degree
at most two. Thus this bare evaluator cannot produce exponent $x_1x_2x_3$
as a formal polynomial. The same holds as a function over a field of size
larger than three, by uniqueness of reduced polynomials of individual
degree below the field size. The script verifies the evaluation-matrix rank
increase when a cubic is added to quadratic functions over $\F_5$.

This toy obstruction excludes auxiliary encodings, LWE ciphertexts, proofs
and randomized representations. It does not rule out bilinear groups as a
component of iO; the supplied construction uses additional structure. It
only prevents an unlicensed iteration of a degree-two pairing from becoming
an assumed multilinear operation.

\paragraph{Stress test.}
A proof supplied by a circuit-generating adversary can support a
proof-sensitive theorem. General iO quantifies over equivalent circuits
without requiring that certificate. Switching to nonuniform reductions
would likewise require an explicit assumption change, not cure the issue
silently. The efficiency improvement in \eref{30006749}{3} retains the proof
hypothesis. Instantiating an ideal oracle is not justified by hybrid
bookkeeping. No toy-model limitation above is claimed for the combined
LWE/LPN/pairing constructions.

\paragraph{Outcome.}
\textbf{Outcome: Exploratory approach with unresolved gap.}
The attempt establishes a polynomial-loss composition criterion, a
quantifier-level failure of polynomial parameter inflation for exponential
hybrids, and a proof-complexity obstruction to universal short certificates.
It constructs no new obfuscator. A successful continuation needs a uniform,
polynomial-loss comparison for arbitrary equivalent circuits without first
requiring polynomially checkable short proofs for all such pairs.

% END RESEARCH ADDITION 142
\subsection{Sources}
\begin{itemize}\raggedright
\item[\textnormal{[S1]}]\hypertarget{source:30006749:1}{} A. Jain, H. Lin, A. Sahai, STOC 2021: 60-73, 2021.
\url{https://doi.org/10.1145/3406325.3451093}.
{\small\textit{Catalogue formulation source.}}
% BEGIN RESEARCH REFERENCES 142
\item[\textnormal{[E1]}]\hypertarget{extra:30006749:1}{} A. Jain, H. Lin and A. Sahai,
\emph{Indistinguishability obfuscation from well-founded assumptions},
arXiv:2008.09317v1, Section 5, Theorems 5.3--5.5; full version underlying [S1].
\url{https://arxiv.org/pdf/2008.09317}.
\item[\textnormal{[E2]}]\hypertarget{extra:30006749:2}{} A. Jain and Z. Jin,
\emph{Indistinguishability obfuscation via mathematical proofs of equivalence},
FOCS 2022, 1023--1034; Cryptology ePrint 2022/1430.
\url{https://eprint.iacr.org/2022/1430}.
\item[\textnormal{[E3]}]\hypertarget{extra:30006749:3}{} Y. Ma, C. Dai and E. Shi,
\emph{Quasi-linear indistinguishability obfuscation via mathematical proofs
of equivalence and applications}, EUROCRYPT 2025, LNCS 15603, 157--186;
ePrint 2025/307, stated equivalence-proof hypothesis.
\url{https://eprint.iacr.org/2025/307}.
\item[\textnormal{[E4]}]\hypertarget{extra:30006749:4}{} A. Jain, H. Lin, J. Luo
and D. Wichs, \emph{The pseudorandom oracle model and ideal obfuscation},
CRYPTO 2023, 233--262; ePrint 2022/1204, abstract and model.
\url{https://eprint.iacr.org/2022/1204}.

% END RESEARCH REFERENCES 142
\end{itemize}

\end{document}
